\documentclass[10pt]{article}
% Shared LaTeX style for MIT 18.06 exam revision notes.
% Compile topic files with Tectonic, XeLaTeX, or LuaLaTeX.

\usepackage[a4paper,margin=0.72in]{geometry}
\usepackage{fontspec}
\usepackage{amsmath,amssymb,mathtools}
\usepackage{booktabs,array,tabularx}
\usepackage{enumitem}
\usepackage{titlesec}
\usepackage{fancyhdr}
\usepackage{microtype}
\usepackage{xcolor}

\setmainfont{Times New Roman}
\setsansfont{Arial}

\setlength{\parindent}{0pt}
\setlength{\parskip}{3.6pt}
\setlength{\abovedisplayskip}{6pt}
\setlength{\belowdisplayskip}{6pt}
\setlength{\abovedisplayshortskip}{4pt}
\setlength{\belowdisplayshortskip}{4pt}
\renewcommand{\arraystretch}{1.08}
\setlength{\tabcolsep}{4.5pt}

\titleformat{\section}
  {\normalfont\large\bfseries}
  {\thesection}
  {0.55em}
  {}
\titlespacing*{\section}{0pt}{10pt}{3pt}

\titleformat{\subsection}
  {\normalfont\normalsize\bfseries}
  {\thesubsection}
  {0.5em}
  {}
\titlespacing*{\subsection}{0pt}{7pt}{2pt}

\setlist[itemize]{leftmargin=1.35em,itemsep=1pt,topsep=2pt,parsep=0pt}
\setlist[enumerate]{leftmargin=1.55em,itemsep=1pt,topsep=2pt,parsep=0pt}

\pagestyle{fancy}
\fancyhf{}
\fancyfoot[C]{\scriptsize\color{black!45}MIT 18.06 Linear Algebra Exam Notes --- \thepage}
\renewcommand{\headrulewidth}{0pt}
\renewcommand{\footrulewidth}{0pt}

\newcommand{\notesubtitle}[1]{%
  {\centering\itshape #1\par}
  \vspace{2pt}
}

\newcommand{\source}[1]{%
  {\centering\small #1\par}
  \vspace{6pt}
}

\newcommand{\labelnote}[2]{%
  \par\smallskip
  \noindent\textbf{#1.}\quad #2
  \par\smallskip
}

\newcommand{\tightlabel}[1]{%
  \par\smallskip
  \noindent\textbf{#1.}\quad
}

\newcolumntype{Y}{>{\raggedright\arraybackslash}X}
\newcolumntype{L}[1]{>{\raggedright\arraybackslash}p{#1}}

\newenvironment{notetable}
  {\small\begin{tabularx}{\linewidth}}
  {\end{tabularx}}


\begin{document}

\begin{center}
{\LARGE 6.3--6.5 Matrix Exponentials, Symmetric Matrices, and Positive Definiteness\par}
\end{center}
\notesubtitle{Concise exam revision notes for continuous systems, symmetric eigenvectors, and positive definite tests.}
\source{Based on handwritten MIT 18.06 revision notes.}

\section{Core Idea}

\labelnote{Core idea}{Eigenvectors decouple matrix problems. For symmetric matrices, the decoupling is especially clean because the eigenvectors can be chosen orthonormal.}

The same eigenvalue idea controls:
\[
\frac{du}{dt}=Au,
\qquad
S=Q\Lambda Q^{T},
\qquad
x^{T}Sx.
\]

\section{Differential Equations}

For the system
\[
\frac{du}{dt}=Au,
\]
if \(A\) has independent eigenvectors \(x_1,\ldots,x_n\), then the solution has the form
\[
u(t)=c_1e^{\lambda_1t}x_1+\cdots+c_ne^{\lambda_nt}x_n.
\]
The constants come from the initial condition:
\[
Sc=u(0),
\qquad
c=S^{-1}u(0).
\]

In matrix form,
\[
u(t)=e^{At}u(0),
\qquad
e^{At}=Se^{\Lambda t}S^{-1}.
\]
Here
\[
e^{\Lambda t}=
\begin{bmatrix}
e^{\lambda_1t}&0&\cdots&0\\
0&e^{\lambda_2t}&\cdots&0\\
\vdots&\vdots&\ddots&\vdots\\
0&0&\cdots&e^{\lambda_nt}
\end{bmatrix}.
\]

\section{Continuous Stability}

\begin{center}
\small
\begin{tabularx}{\linewidth}{@{}L{0.32\linewidth}Y@{}}
\toprule
Eigenvalues of \(A\) & Behavior of \(e^{At}\) \\
\midrule
All \(\operatorname{Re}\lambda<0\) & Decays to zero. \\
Some \(\operatorname{Re}\lambda>0\) & Blows up in that direction. \\
One \(\lambda=0\), rest \(\operatorname{Re}\lambda<0\) & Possible steady state. \\
Complex \(\lambda=a+ib\) & \(a\) controls growth/decay; \(b\) controls oscillation. \\
\bottomrule
\end{tabularx}
\end{center}

\labelnote{Exam memory}{For powers use \(|\lambda|\). For differential equations use \(\operatorname{Re}\lambda\).}

\section{Defective Matrices}

A repeated eigenvalue may not give enough eigenvectors. Then \(S\) is not invertible, so
\[
A=S\Lambda S^{-1}
\]
is impossible.

For a defective repeated eigenvalue, a generalized eigenvector can replace a missing eigenvector:
\[
(A-\lambda I)w=x.
\]
This creates terms involving \(te^{\lambda t}\) in differential equations.

\labelnote{Scope}{For exam revision, first try ordinary eigenvectors. Use generalized eigenvectors only when the matrix is defective.}

\section{Symmetric Matrices}

For a real symmetric matrix
\[
S=S^{T}.
\]
The key facts are:
\begin{itemize}
  \item all eigenvalues are real;
  \item eigenvectors for different eigenvalues are orthogonal;
  \item the eigenvectors can be chosen orthonormal.
\end{itemize}

Therefore a symmetric matrix factors as
\[
S=Q\Lambda Q^{T}.
\]
Equivalently,
\[
S=\lambda_1q_1q_1^{T}+\cdots+\lambda_nq_nq_n^{T}.
\]

\labelnote{Exam trigger}{If the matrix is symmetric, look for orthogonal eigenvectors and use \(Q^{T}=Q^{-1}\).}

\section{Positive Definite Matrices}

For a symmetric matrix \(S\), positive definite means
\[
x^{T}Sx>0
\qquad
\text{for every }x\ne0.
\]
Equivalent tests:

\begin{center}
\small
\begin{tabularx}{\linewidth}{@{}L{0.30\linewidth}Y@{}}
\toprule
Test & Positive definite condition \\
\midrule
Eigenvalues & All \(\lambda_i>0\). \\
Pivots & All pivots are positive. \\
Leading determinants & All leading principal determinants are positive. \\
Quadratic form & \(x^{T}Sx>0\) for all \(x\ne0\). \\
Factorization & \(S=R^{T}R\) with independent columns in \(R\). \\
\bottomrule
\end{tabularx}
\end{center}

For \(2\times2\)
\[
S=
\begin{bmatrix}
a&b\\
b&c
\end{bmatrix},
\]
the quick test is
\[
a>0,
\qquad
ac-b^{2}>0.
\]

\section{Positive Semidefinite}

Positive semidefinite means
\[
x^{T}Sx\ge0
\qquad
\text{for every }x.
\]
Then eigenvalues are nonnegative:
\[
\lambda_i\ge0.
\]
It may be singular:
\[
\det S=0
\]
is allowed.

\labelnote{Common mistake}{Positive semidefinite is not the same as positive definite. Zero eigenvalues are allowed only in the semidefinite case.}

\section{Similarity and Transpose}

Similar matrices have the form
\[
B=M^{-1}AM.
\]
They have the same eigenvalues because they represent the same linear map in different coordinates.

A matrix and its transpose have the same eigenvalues:
\[
\det(A^{T}-\lambda I)=\det(A-\lambda I).
\]
But their eigenvectors do not have to be the same.

\section{Quadratic Forms and Axes}

A symmetric matrix rotates quadratic forms into principal axes:
\[
x^{T}Sx=(Q^{T}x)^{T}\Lambda(Q^{T}x).
\]
For an ellipse, the eigenvectors give the axis directions and the eigenvalues control the axis lengths.

\section{Common Mistakes}

\begin{center}
\small
\begin{tabularx}{\linewidth}{@{}L{0.40\linewidth}Y@{}}
\toprule
Mistake & Correct exam habit \\
\midrule
Using \(|\lambda|\) for differential equation stability. & Use \(\operatorname{Re}\lambda\). \\
Assuming repeated eigenvalues give enough eigenvectors. & Check the eigenspace dimension. \\
Using \(Q\Lambda Q^{T}\) for a non-symmetric matrix. & Reserve that clean form for symmetric matrices. \\
Testing positive definiteness without symmetry. & The standard tests here assume a symmetric matrix. \\
Confusing semidefinite with definite. & Positive definite needs strict positivity and no zero eigenvalues. \\
Assuming \(A\) and \(A^{T}\) have the same eigenvectors. & They have the same eigenvalues, not necessarily the same eigenvectors. \\
\bottomrule
\end{tabularx}
\end{center}

\section{Final Exam Checklist}

\tightlabel{Checklist}
\begin{enumerate}
  \item For \(u'=Au\), find eigenvalues and eigenvectors of \(A\).
  \item Write \(u(t)=\sum c_ie^{\lambda_it}x_i\) when \(A\) has enough eigenvectors.
  \item Use \(Sc=u(0)\) to find constants.
  \item Use \(\operatorname{Re}\lambda\) to decide continuous stability.
  \item Check symmetry before using \(S=Q\Lambda Q^{T}\).
  \item For a symmetric matrix, choose eigenvectors orthonormal.
  \item Test positive definiteness by eigenvalues, pivots, leading determinants, or \(x^{T}Sx\).
  \item Treat zero eigenvalues as semidefinite, not definite.
  \item Use similarity to preserve eigenvalues, not necessarily eigenvectors.
\end{enumerate}

\end{document}
